\documentclass[11pt]{amsart}

\usepackage[T1]{fontenc}
\usepackage{amsmath,amssymb}
\usepackage{microtype}
\usepackage[
  colorlinks=true,
  linkcolor=blue,
  citecolor=blue,
  urlcolor=blue,
  pdftitle={DRAFT: Degree, defect and separators in K7-minus-minor-free critical graphs},
  pdfauthor={Cavit Erginsoy}
]{hyperref}
\setlength{\emergencystretch}{1em}

\newtheorem{theorem}{Theorem}[section]
\newtheorem{lemma}[theorem]{Lemma}
\newtheorem{corollary}[theorem]{Corollary}
\newtheorem{proposition}[theorem]{Proposition}
\newcommand{\minor}{\preccurlyeq}
\newcommand{\N}{N}

\title[DRAFT: Degree, defect and separators]
{\textbf{DRAFT}\\[0.5ex]
Degree, defect and separators in
$K_7^{-}$-minor-free critical graphs}
\author{Cavit Erginsoy}
\date{Draft, 13 September 2026}

\subjclass[2020]{05C15, 05C40}
\keywords{Hadwiger's conjecture, graph minors, contraction-critical graphs,
rooted minors}

\begin{document}

\begin{abstract}
Let $G$ be a minor-minimal non-six-colourable graph with no $K_7^{-}$
minor.  For each $i$, let $n_i$ count its degree-$i$ vertices, and put
$\tau=\sum_{i\geq10}(i-9)n_i$.  We prove that $G$ is seven-connected,
has minimum degree eight, contains no $K_5$ subgraph, and satisfies
$n_8\geq27+\tau$.  Every degree-eight vertex is incident with an edge
lying in at most three triangles, and deleting any seven-vertex cut
leaves exactly two components.  The bound on $n_8$ follows from the broader
statement that every $r$-connected $K_7^{-}$-minor-free graph with $r\geq6$
and at least four times as many edges as vertices has defect
$9|V|-2|E|\geq20+r$.  This statement uses a computer-assisted finite
lemma on nine- and ten-vertex quotients.  We also give computation-free
proofs of $n_8\geq26+\tau$, the incident-edge conclusion and the
separator conclusion.
\end{abstract}

\maketitle
\markboth{DRAFT: DEGREE, DEFECT AND SEPARATORS}{CAVIT ERGINSOY: DRAFT}

\section{Introduction}

Hadwiger's conjecture states that every graph of chromatic number $t$
contains a $K_t$ minor~\cite{Hadwiger1943}.  It is known for $t\leq6$, with
the last case due to Robertson, Seymour and Thomas
\cite{RobertsonSeymourThomas1993}, and remains open for $t\geq7$.  Write
$K_t^{-}$ for the graph obtained from $K_t$ by deleting one edge, and
$K_7^{\vee}$ for the graph obtained from $K_7$ by deleting two incident
edges.  Albar proved that every $K_7^{-}$-minor-free graph is
seven-colourable~\cite[Theorem~1]{Albar2014}.  Norin and Totschnig proved
that every $K_7^{\vee}$-minor-free graph is
six-colourable~\cite[Theorem~4]{NorinTotschnig2025}.  They asked whether
$K_7^{-}$ may replace $K_7^{\vee}$ in that statement
\cite[Conjecture~21]{NorinTotschnig2025}.

A graph is \emph{minor-minimal non-six-colourable} if it is not
six-colourable and every proper minor is six-colourable.  The following
theorem describes such a graph under the exclusion in Conjecture 21.

\begin{theorem}\label{thm:main}
Let $G$ be a minor-minimal non-six-colourable $K_7^{-}$-minor-free graph.
For each $i$, let $n_i$ denote the number of vertices of degree $i$, and put
$\tau=\sum_{i\geq10}(i-9)n_i$.
Then $G$ is seven-connected and
\[
        \delta(G)\geq8,
        \qquad |E(G)|\geq4|V(G)|.
\]
Moreover, $G$ contains no $K_5$ subgraph,
\[
        n_8\geq27+\tau,
\]
and the neighbourhood of every degree-eight vertex is $K_4$-free and has
independence number three.  Every such vertex is incident with an edge
lying in at most three triangles.  Every seven-vertex cut of $G$ leaves
exactly two components.

The bound $n_8\geq27+\tau$ uses the finite lemma
\ref{lem:finite-quotients}.  All the other conclusions, together with the
weaker bound $n_8\geq26+\tau$, have computation-free proofs.
\end{theorem}

These are necessary conditions for a counterexample to Conjecture 21.
A sufficient open assertion
is that every seven-connected graph $H$ satisfying
$|E(H)|\geq4|V(H)|$ contains a $K_7^{-}$ minor.  The bound on $n_8$
follows from the next theorem.

\begin{theorem}\label{thm:defect-ladder}
Let $r\geq6$, and let $H$ be an $r$-connected $K_7^{-}$-minor-free graph
with $|E(H)|\geq4|V(H)|$.  Then
\[
                    9|V(H)|-2|E(H)|\geq20+r.
\]
The proof uses the finite lemma \ref{lem:finite-quotients}.
\end{theorem}

The proofs use linked cliques and rooted models to exclude vertices of
degree seven and $K_5$ subgraphs.  Section~\ref{sec:degree-eight} then gives
the computation-free bound $n_8\geq26+\tau$.  Theorem~\ref{thm:defect-ladder}
uses a density-preserving contraction and induction on connectivity
(Section~\ref{sec:defect}).  The separator proof aligns boundary colourings
obtained from proper minors and extends precolourings over planar graphs
(Section~\ref{sec:seven-cuts}).

Mayer obtained a finite catalogue of candidate neighbourhood graphs for
degree-seven vertices in six-contraction-critical graphs~\cite{Mayer1993}.
In the minor-minimal $K_7^{\vee}$-minor-free setting, Norin and Totschnig
proved that every degree-seven vertex belongs to a $K_5$ contained in its
closed neighbourhood~\cite[Claim~4.4]{NorinTotschnig2025}.  We first
determine the neighbourhood under $K_7^{-}$ exclusion, then rule out degree
seven.

Rolek and Song proved the next result when $r=6$
~\cite[Lemma~1.9]{RolekSong2017}.  We shall use the uniform form with
$r=5$.

\begin{theorem}\label{thm:two-cliques}
For every $r\geq2$, if $G$ is an $(r+1)$-connected graph containing two
distinct $K_r$ subgraphs, then $K_{r+2}^{-}\minor G$.
\end{theorem}

The connectivity hypothesis is sharp.  The graph $K_{r+1}$ is
$r$-connected and contains two distinct $K_r$ subgraphs, but has no
$K_{r+2}^{-}$ minor.

The following rooted closure will exclude the remaining degree-seven
neighbourhood.

\begin{theorem}\label{thm:rooted-closure}
Let $G$ be a six-connected graph of order $n$, let $v\in V(G)$ have degree
$d\geq5$, and suppose that $G[\N(v)]$ contains a $K_4$ subgraph.  If
\[
        |E(G)|\geq4n+d-13,
\]
then $K_7^{-}\minor G$.
\end{theorem}

A \emph{minor model} of a graph $H$ in $G$ is a family
$(B_x:x\in V(H))$ of pairwise disjoint connected vertex sets such that
$B_x$ and $B_y$ are adjacent whenever $xy\in E(H)$.  A model is
\emph{rooted} at prescribed vertices when its branch sets contain them
one-to-one.  We write $H\minor G$ when $H$ is a minor of $G$.  All graphs
are finite and simple.

\section{Two linked cliques}

\begin{lemma}\label{lem:disjoint-cliques}
Let $r\geq2$.  If an $(r+1)$-connected graph $H$ contains two
vertex-disjoint $K_r$ subgraphs, then $K_{r+2}^{-}\minor H$.  The branch
sets can be chosen to meet the union of the two cliques.
\end{lemma}

\begin{proof}
Write the two cliques as
\[
        A=\{a_1,\ldots,a_r\},
        \qquad B=\{b_1,\ldots,b_r\}.
\]
By Menger's theorem there are $r$ pairwise vertex-disjoint $A$ to $B$
paths whose open interiors avoid $A\cup B$.  Relabel them as paths $P_i$
from $a_i$ to $b_i$, and put $U=\bigcup_iV(P_i)$.

Call a path $Q$ a \emph{usable cross-path} if its ends lie on distinct
paths $P_i,P_j$, its open interior avoids $U$, and its ends are not both
$A$-ends or both $B$-ends.  Suppose first that $Q$ exists.  Split each of
$P_i$ and $P_j$ across an edge into a nonempty $A$-prefix and a nonempty
$B$-suffix, choosing the split so that the ends of $Q$ belong to pieces of
opposite kinds.  Absorb the open interior of $Q$ into the piece containing
one end.  The four resulting pieces are connected and disjoint.  The two
$A$-pieces are adjacent through $A$, the two $B$-pieces are adjacent
through $B$, the two pieces from each original path are adjacent across
its splitting edge, and $Q$ supplies one of the two remaining
cross-adjacencies.  Thus at most one adjacency among the four pieces is
missing.  Each unsplit path is adjacent to every split piece through its
two clique ends, and the unsplit paths are pairwise adjacent.  These
$r+2$ sets form a $K_{r+2}^{-}$ model.

Suppose now that no usable cross-path exists.  Let $C$ be a component of
$H-U$.  If $C$ has neighbours on two distinct paths, a path through $C$
between two such attachments is unusable only when both attachments are
$A$-ends or both are $B$-ends.  Comparing all pairs of attachments shows
that all neighbours of $C$ are $A$-ends or all are $B$-ends.  Hence
$|\N_H(C)|\leq r$, contrary to $(r+1)$-connectivity.  Every component of
$H-U$ therefore attaches to only one $P_i$.

For each $i$, let $W_i$ consist of the open interior of $P_i$ and the
components of $H-U$ attached to $P_i$.  There is no edge from $W_i$ to
$P_j$ for $j\ne i$: an edge from the path interior would be a usable
cross-path, and an edge from a component would contradict its unique
attachment.  Hence
\[
        \N_H(W_i)\subseteq\{a_i,b_i\}.
\]
Since $r\geq2$, a nonempty $W_i$ would give a vertex cut of order at most
two.  Thus every $W_i$ is empty.  It follows that $H$ consists of $A$, $B$
and the matching edges $a_ib_i$: any other edge between $A$ and $B$ would
be usable.  This graph has minimum degree $r$, again contradicting
$(r+1)$-connectivity.
\end{proof}

\begin{proof}[Proof of Theorem~\ref{thm:two-cliques}]
Let $L_1,L_2$ be distinct $K_r$ subgraphs and put
$s=|V(L_1)\cap V(L_2)|$.  If $s\leq r-2$, delete the common set $Z$.
The graph $G-Z$ is $(r+1-s)$-connected and contains the two disjoint
$K_{r-s}$ subgraphs $L_1-Z$ and $L_2-Z$.  Lemma~\ref{lem:disjoint-cliques}
gives a $K_{r-s+2}^{-}$ model in which every branch set contains a vertex
of $(L_1\cup L_2)-Z$.  The vertices of $Z$, taken as singleton branch
sets, are complete to this model and to one another.  They complete a
$K_{r+2}^{-}$ model.

It remains that $s=r-1$.  The union $X=V(L_1)\cup V(L_2)$ induces a graph
containing $K_{r+1}^{-}$.  Since an $(r+1)$-connected graph has at least
$r+2$ vertices, $G-X$ is nonempty.  Let $C$ be a component of $G-X$.  If
$C$ missed a vertex of $X$, then $\N_G(C)$ would be a cut of order at most
$r$.  Thus $C$ is adjacent to every vertex of $X$.  The vertices of $X$
as singleton branch sets, together with $C$, form a $K_{r+2}^{-}$ model.
\end{proof}

\begin{corollary}\label{cor:unique-k5}
Every six-connected $K_7^{-}$-minor-free graph contains at most one $K_5$
subgraph.
\end{corollary}

\section{Degree-seven neighbourhoods}

A graph is \emph{$k$-contraction-critical} if it has chromatic number $k$
and every proper minor is $(k-1)$-colourable.  We use two standard facts.
A noncomplete $k$-contraction-critical graph is seven-connected when
$k\geq7$~\cite[Satz~6]{Mader1968Connectivity}; see also
\cite[Theorem~1.8]{RolekSong2017}.  If $G$ is $k$-contraction-critical and
$v\in V(G)$, then
\[
        \alpha(G[\N(v)])\leq d(v)-k+2.                 \tag{3.1}
\]
This inequality goes back to Dirac~\cite{Dirac1960}; we use the formulation
in \cite[Lemma~1.6(i)]{RolekSong2017}.

We also use the following specialisation of Theorem~7 and property $(*)$
of Kriesell and Mohr~\cite[Theorem~7]{KriesellMohr2019}.  Let a graph $J$
be properly coloured with five colours, and choose one root from each
colour class.  Let $D$ be a graph on the five roots with at most six edges.
If the ends of every edge $xy\in E(D)$ lie in the same bichromatic
component of $J$ on their two colours, then $J$ contains pairwise disjoint
connected sets, one containing each root, that are adjacent for every edge
of $D$.

\begin{lemma}\label{lem:rooted-k5}
Let $G$ be a seven-chromatic graph such that every proper minor is
six-colourable and $K_7^{-}\not\minor G$.  Let $v$ have degree seven, put
$H=G[\N(v)]$, and let $ab$ be a nonedge of $H$.  If
\[
        R=\N(v)\setminus\{a,b\},
\]
then $G-v-\{a,b\}$ contains a $K_5$ minor rooted at $R$.
\end{lemma}

\begin{proof}
Contract the two edges $va$ and $vb$ to one vertex.  A six-colouring of
the resulting proper minor induces a six-colouring $\varphi$ of $G-v$ in
which $a$ and $b$ have one common colour, say $0$.  Every vertex of $R$
avoids colour $0$.  The five vertices of $R$ use the other five colours
distinctly, since a colour absent from $\N(v)$ could be assigned to $v$.

Put $F=\overline H$.  Inequality~(3.1) gives $\alpha(H)\leq2$, so $F$ is
triangle-free.  If $xy\in E(F[R])$ and $x,y$ have colours $i,j$, then
$x,y$ lie in the same component of the subgraph of $G-v$ induced by
colours $i,j$.  Otherwise, interchanging $i$ and $j$ on the component
containing $x$ would remove colour $i$ from $\N(v)$, allowing that colour
at $v$.

Delete the colour-$0$ class from $G-v$.  The resulting graph is properly
five-coloured, with $R$ a transversal of its colour classes.  By Mantel's
theorem, the triangle-free graph $F[R]$ has at most six edges.  Apply the
Kriesell and Mohr result with demand graph $F[R]$.  It gives five disjoint
connected sets rooted at $R$ and adjacent for every nonedge of $H[R]$.
Every other pair is joined by its edge in $H[R]$, so the five sets form the
required rooted $K_5$ model.
\end{proof}

\begin{theorem}\label{thm:local-classification}
Under the hypotheses of Lemma~\ref{lem:rooted-k5}, every degree-seven
vertex $v$ has $G[\N(v)]\cong K_4\mathbin{\dot\cup}K_3$ and hence
belongs to a $K_5$ subgraph.
\end{theorem}

\begin{proof}
Put $H=G[\N(v)]$ and $F=\overline H$.  By (3.1), $F$ is triangle-free.
It is nonempty, since otherwise $G[\N[v]]$ would contain $K_7^{-}$.

Every nonisolated vertex of $F$ has degree at least three.  To see this,
choose an edge $ab\in E(F)$ and use the rooted model of
Lemma~\ref{lem:rooted-k5} on $R=\N(v)\setminus\{a,b\}$.  If $d_F(a)=1$,
then $a$ is adjacent in $H$ to every root; the rooted model, $\{a\}$ and
$\{v\}$ form a $K_7$ model.  If $d_F(a)=2$, with other neighbour $x$,
then triangle-freeness gives $bx\in E(H)$.  Absorb $b$ into the branch set
rooted at $x$.  Together with $\{a\}$ and $\{v\}$, the resulting seven
branch sets have at most one missing adjacency, again a contradiction.

Every nontrivial component of $F$ therefore has minimum degree at least
three and hence at least six vertices.  Since $|V(F)|=7$, there is one
nontrivial component and at most one isolated vertex.  If the component
has six vertices, Mantel's theorem and its equality case give
$F\cong K_{3,3}\mathbin{\dot\cup}K_1$.

Suppose instead that the component has seven vertices.  No vertex has
degree at least five: its independent neighbourhood would contain a
vertex of degree at most two.  Thus every degree is three or four.  The
degree sum is even, so some vertex $z$ has degree four.  Its four
neighbours are independent.  Each is adjacent to both vertices outside
$\N_F[z]$ to have degree at least three, and those two outside vertices
are nonadjacent by triangle-freeness.  Hence $F\cong K_{3,4}$.  Taking
complements leaves $K_4\mathbin{\dot\cup}K_3$ or
$K_1\vee(K_3\mathbin{\dot\cup}K_3)$ for $H$.  Mader's theorem makes $G$
seven-connected, so Corollary~\ref{cor:unique-k5} gives at most one $K_5$
subgraph.  The second possibility would place $v$ in two distinct $K_5$
subgraphs, leaving the claimed neighbourhood.
\end{proof}

\section[Rooted model closure]{Rooted
\texorpdfstring{$K^*_{4,2}$}{K-star(4,2)} closure}

For a graph $F$ and $S\subseteq V(F)$, say that $(F,S)$ is
\emph{internally $k$-connected} if there is no separation $(A,B)$ of $F$
with $S\subseteq A$, $B\setminus A\ne\varnothing$, and
$|A\cap B|<k$.  For a four-set $Z$, a $Z$-rooted $K^*_{4,2}$ model
consists of six pairwise disjoint connected sets: four rooted at $Z$ and
two others $U,V$, each adjacent to the other five.  Adjacencies among the
rooted sets are not required.

\begin{theorem}[Norin and Totschnig {\cite[Lemma~12]{NorinTotschnig2025}}]
\label{thm:rooted-two-helper-bound}
Let $F$ be a graph and let $Z\subseteq V(F)$ have order four.  If $(F,Z)$
is internally four-connected and $F$ has no $Z$-rooted $K^*_{4,2}$ model,
then
\[
        |E(F)|\leq4|V(F)|-10.                          \tag{4.1}
\]
\end{theorem}

The following placement lemma puts one prescribed vertex in a non-root
branch set.

\begin{lemma}\label{lem:fifth-root-helper}
Let $F$ be a graph and let $S=Z\mathbin{\dot\cup}\{x\}$, where $|Z|=4$.
Suppose that $(F,S)$ is internally five-connected.  If $F$ has a
$Z$-rooted $K^*_{4,2}$ model, then it has one in which $x$ belongs to a
non-root branch set.
\end{lemma}

\begin{proof}
Choose a $Z$-rooted model that maximises $|U|+|V|$ over its two non-root
branch sets $U,V$ and, subject to that, minimises
$\sum_{i=1}^4|R_i|$ over its rooted branch sets $R_1,\ldots,R_4$.  Write
$z_i$ for the root in $R_i$, and put
\[
        P_i=\{r\in R_i:r\text{ has a neighbour in }U\cup V\}.
\]
We claim that $|P_i|=1$.  Both $U$ and $V$ are adjacent to $R_i$, so $P_i$ is
nonempty.  If it had at least two vertices, there would be distinct
$u,w\in R_i$ such that $u$ has a neighbour in $U$ and $w$ has a neighbour in $V$; otherwise the two
contact sets would be the same singleton.  Take a minimal tree in
$F[R_i]$ containing $z_i,u,w$.  By the minimal choice of the rooted branch
sets, this tree spans $R_i$.  One of $u,w$, say $u$, is a leaf different
from $z_i$.  Move $u$ from $R_i$ into $U$.  Both altered branch sets remain
connected.  The former tree edge from $u$ to $R_i\setminus\{u\}$ preserves
the contact between $R_i$ and $U$, while $w$ preserves the contact between
$R_i$ and $V$.  Every other required adjacency is unchanged, contrary to
the choice of the model.  Hence $|P_i|=1$.

No component outside the six branch sets can have a neighbour in
$U\cup V$, since it could be absorbed into the adjacent branch set.  It follows that
\[
        |N_F(U\cup V)\setminus(U\cup V)|\leq4,          \tag{4.2}
\]
with at most one external neighbour in each rooted branch set.  If
$x\notin U\cup V$, (4.2) would give a separation of $(F,S)$ of order at
most four whose second open side is the nonempty set $U\cup V$.  This
contradicts internal five-connectivity.  Thus $x\in U\cup V$.
\end{proof}

\begin{proof}[Proof of Theorem~\ref{thm:rooted-closure}]
Choose a four-set $Z\subseteq N_G(v)$ inducing $K_4$, choose
$x\in N_G(v)\setminus Z$, and put $F=G-v$.  Deleting one vertex from a
six-connected graph leaves a five-connected graph.  In particular,
$(F,Z)$ is internally four-connected and $(F,Z\cup\{x\})$ is internally
five-connected.  Moreover,
\[
 \begin{aligned}
 |E(F)|&=|E(G)|-d\\
       &\geq4n-13
        =4|V(F)|-9.
 \end{aligned}                                          \tag{4.3}
\]
Theorem~\ref{thm:rooted-two-helper-bound} supplies a $Z$-rooted
$K^*_{4,2}$ model in $F$, and Lemma~\ref{lem:fifth-root-helper} lets us
choose it with $x$ in one non-root branch set.

The four rooted branch sets are pairwise adjacent through the edges of
$G[Z]=K_4$.  Together with the two non-root branch sets they form a $K_6$
model.  The singleton branch set $\{v\}$ is adjacent to all four rooted
branch sets and to the one containing $x$; it may miss only the other.
These seven branch sets form a $K_7^{-}$ model in $G$.
\end{proof}

\section{Consequences for a critical graph}

Jakobsen proved that an $n$-vertex graph with $n\geq7$ and at least
$\frac92n-12$ edges contains a $K_7^{-}$ minor or is a
$(K_{2,2,2,2},K_6,4)$-cockade~\cite[Theorem~4]{Jakobsen1983}; see also
\cite[Theorem~2]{Albar2014}.

Write $D(G)=9|V(G)|-2|E(G)|$ for the \emph{defect} of $G$.
The \emph{codegree} $c(xy)=|N(x)\cap N(y)|$ counts the triangles containing
an edge $xy$.  Contracting an edge reduces connectivity by at most one:
a cut avoiding the contraction vertex lifts unchanged, while a cut
containing it lifts after that vertex is replaced by the two ends.
For an edge $e$ of codegree $c$,
\begin{equation}\label{eq:contraction}
 |E(G/e)|=|E(G)|-1-c,\qquad D(G/e)=D(G)-7+2c.             \tag{5.1}
\end{equation}

\begin{lemma}\label{lem:jakobsen-defect}
Every five-connected $K_7^{-}$-minor-free graph $G$ of order at least nine
satisfies $D(G)\geq25$.
\end{lemma}

\begin{proof}
The order excludes the two base cockades $K_6$ and $K_{2,2,2,2}$.
A nontrivial cockade has a four-vertex cut, contrary to five-connectivity.
Jakobsen's strict inequality $2|E(G)|<9|V(G)|-24$ and integrality give
$D(G)\geq25$.
\end{proof}

\begin{proposition}\label{prop:critical-base}
Every graph satisfying the hypotheses of Theorem~\ref{thm:main} is
seven-connected, has minimum
degree eight, has no $K_5$ subgraph, and satisfies
$|E(G)|\geq4|V(G)|$ and $n_8\geq25+\tau$.
\end{proposition}

\begin{proof}
Albar's theorem gives $\chi(G)=7$, and Mader's theorem gives
seven-connectivity.  By Theorem~\ref{thm:local-classification} and
Corollary~\ref{cor:unique-k5}, every degree-seven vertex lies in the unique
possible $K_5$, so $n_7\leq5$.  Degree summation gives
$2|E(G)|\geq8|V(G)|-5$ and hence $|E(G)|\geq4|V(G)|-2$.
If $v$ had degree seven, its neighbourhood would contain a $K_4$ and this
density would exceed the threshold $4|V(G)|-6$ in
Theorem~\ref{thm:rooted-closure}.  Thus no such vertex exists:
$\delta(G)\geq8$ and $|E(G)|\geq4|V(G)|$.

Put $q=|E(G)|-4|V(G)|$, so $q\geq0$.  Suppose that a $K_5$ subgraph has
vertex set $K$.  For each $w\in K$, the neighbourhood of $w$ contains the
$K_4$ on $K\setminus\{w\}$.  Since $G$ has no $K_7^{-}$ minor, the
contrapositive of Theorem~\ref{thm:rooted-closure} gives
\[
        4|V(G)|+q<4|V(G)|+d_G(w)-13,
\]
and hence $d_G(w)\geq q+14$.  The five vertices of $K$ would contribute at
least $5(q+6)$ to
\[
        \sum_{z\in V(G)}(d_G(z)-8)=2q,
\]
which is impossible.  Thus $G$ contains no $K_5$ subgraph.

The minimum degree ensures $|V(G)|\geq9$.  Lemma~\ref{lem:jakobsen-defect}
and degree summation now give
\[
        25\leq9|V(G)|-2|E(G)|=n_8-\tau.
\]
Therefore $n_8\geq25+\tau$.
\end{proof}

\section{Degree-eight vertices without computation}\label{sec:degree-eight}

\begin{lemma}\label{lem:independent-triple}
Let $G$ satisfy the hypotheses of Theorem~\ref{thm:main}.  Every
degree-eight vertex $v$ satisfies $\alpha(G[N(v)])=3$.
\end{lemma}

\begin{proof}
Dirac's inequality (3.1) gives the upper bound three.  Put $J=G[N(v)]$.
By Proposition~\ref{prop:critical-base}, $J$ has no $K_4$ subgraph and
$|V(G)|\geq n_8\geq25$, so $G-N[v]$ is nonempty.  If $\alpha(J)\leq2$,
Rolek, Song and Thomas~\cite[Lemma~2.1]{RolekSongThomas2022} show that $J$
contains a spanning square of an eight-cycle.  Label its vertices
cyclically $0,\ldots,7$.  An exterior component has at least seven
neighbours in $N(v)$, by seven-connectivity.  Contract that component to
$c$ and delete the other exterior vertices.  By cyclic symmetry and edge
deletion, take $c$ adjacent to every cycle vertex except $0$.  The sets
\[
       \{0,7,2\},\quad\{3\},\quad\{4\},\quad\{1,v\},
       \quad\{6\},\quad\{5\},\quad\{c\}
\]
are connected and pairwise adjacent except for $\{3\}:\{6\}$.  They
form a forbidden $K_7^{-}$ model.  Hence $\alpha(J)=3$.
\end{proof}

\begin{lemma}\label{lem:good-deletion}
Let $J$ be an eight-vertex graph with minimum degree at least four, no
$K_4$ subgraph, and independence number three.  Call $r$ good if $J-r$
has a $K_5^{-}$ minor.  Every vertex of $J$ has a good neighbour.
\end{lemma}

\begin{proof}
The graph $J$ is three-connected.  Indeed, a component has order at least
five, and a component behind a cutvertex has order at least four, excluding
both kinds of separation on eight vertices.  A two-cut would leave two
three-vertex components, each a triangle complete to the cut.  An edge in
the cut gives a $K_4$; an independent cut gives independence number two.
Both are impossible.

Suppose first that $J$ is four-connected.  Then $J-r$ is three-connected
for every $r$.  Wood and Woodall~\cite[Lemma~4.2.1]{WoodWoodall2009}
classify three-connected $K_5^{-}$-minor-free graphs as wheels, the
triangular prism and $K_{3,3}$.  If some $r$ is not good, order seven
forces $J-r$ to be a wheel with six rim vertices.  Every rim vertex must
see $r$ to have degree at least four, and the hub cannot see $r$ because
$J$ has no $K_4$.  Thus $J=\overline{K_2}\vee C_6$.  With apices $p,q$
and rim $v_0,\ldots,v_5$, deletion of $v_2$ leaves the $K_5^{-}$ model
\[
       \{p,v_0\},\quad\{q,v_1\},\quad\{v_3\},
       \quad\{v_4\},\quad\{v_5\},
\]
missing only $\{v_3\}:\{v_5\}$.  Every rim vertex is therefore good,
and the rim vertices meet every neighbourhood.  If no bad $r$ exists,
the conclusion is immediate.

Now let $S$ be a three-cut.  Its two components have orders two and three,
since each has order at least two.  Write the first as $A=\{a,a'\}$ and
the second as $B$.  The edge $aa'$ is present and both ends are complete
to $S$.  Consequently $S$ is independent.  Every vertex of $B$ has a
neighbour in $B$, so $J[B]$ is a path or triangle.

If $B=b_1b_2b_3$ is a path, its ends are complete to $S$ and $b_2$ has at
least two neighbours in $S$.  The following models establish that all
vertices in $A\cup B$ are good; symmetry handles $a'$ and $b_3$.
In the table, $b_1s_1$ abbreviates the branch set $\{b_1,s_1\}$, and
similarly for the other entries.
\[
\begin{array}{c|c|c}
\text{deleted vertex}&\text{five branch sets}&\text{only missing pair}\\\hline
a&b_1s_1,b_3s_2,b_2,a',s_3&b_2:a'\\
b_1&as_1,b_3s_2,b_2,a',s_3&b_2:a'\\
b_2&as_1,b_1s_2,b_3,a',s_3&b_3:a'
\end{array}
\]
Here $S=\{s_1,s_2,s_3\}$ is labelled with $b_2s_3$ present in the first
row, and with $b_2s_1,b_2s_3$ present in the second.

If $B$ is a triangle, each of its vertices sees at least two vertices of
$S$, and no vertex of $S$ sees all of $B$.  Its missing incidences with
$S$ therefore form a perfect matching; label these $b_is_i$.
After deleting $a$ use
\[
       s_1b_2,\quad s_2a',\quad b_1,\quad s_3,\quad b_3,
\]
which misses only $s_3:b_3$; after deleting $b_1$ use
\[
       s_1b_2b_3,\quad s_2,\quad s_3,\quad a,\quad a',
\]
which misses only $s_2:s_3$.  Symmetry again makes all vertices in
$A\cup B$ good.  In either case, the vertices of $A$ see one another,
every vertex of $B$ has a neighbour in $B$, and every vertex of $S$ sees
$A$.  Thus every vertex has a good neighbour.
\end{proof}

\begin{proposition}\label{prop:critical-three-edge}
Let $G$ satisfy the hypotheses of Theorem~\ref{thm:main}.  Every
degree-eight vertex is incident with an edge $e$ satisfying $c(e)\leq3$,
and $n_8\geq26+\tau$.  Both conclusions have computation-free proofs.
\end{proposition}

\begin{proof}
Fix a degree-eight vertex $v$ and put $J=G[N(v)]$.  By
Proposition~\ref{prop:critical-base} and Lemma~\ref{lem:independent-triple},
$J$ is $K_4$-free and has independence number three.  If every incident
edge has codegree at least four, then $\delta(J)\geq4$.  An exterior
component exists because $|V(G)|\geq25$, and seven-connectivity gives at
least seven boundary neighbours.  Contract it to $c$ and delete the other
exterior components.  Let $x$ be the vertex of $J$ missed by $c$, or any
vertex if none is missed.  By Lemma~\ref{lem:good-deletion}, $x$ has a
neighbour $r$ such that $J-r$ has a $K_5^{-}$ model $M_1,\ldots,M_5$.
Then $cr$ is an edge, and
\[
                  \{v\},\quad\{c,r\},\quad M_1,\ldots,M_5
\]
form a $K_7^{-}$ model: the first two sets are adjacent through $vr$ and
are adjacent to every $M_i$.  The sole possible missing contact from $c$
is a singleton $M_i=\{x\}$, supplied instead by $rx$.  This contradiction
proves that some edge $e=vx$ has $c(e)=d_J(x)\leq3$.

The quotient $G/e$ is six-connected, has no $K_7^{-}$ minor and has order
at least twenty-four.  Lemma~\ref{lem:jakobsen-defect} gives $D(G/e)\geq25$,
while \eqref{eq:contraction} gives $D(G)\geq D(G/e)+1\geq26$.
The identity $D(G)=n_8-\tau$ proves the bound.
\end{proof}

\section{Connectivity and defect}\label{sec:defect}

\begin{lemma}\label{lem:eight-complement}
Every eight-vertex graph $J$ of minimum degree at least five contains a
$K_6^{-}$ minor.
\end{lemma}

\begin{proof}
The complement of $J$ has maximum degree at most two.  Extend it to an
edge-maximal graph $F$ with that property.  Every component of $F$ is a
cycle, except possibly one component which is $K_1$ or $K_2$: any longer
path could be closed, and deficient vertices in two components could be
joined.  The seven possibilities for $F$ and explicit models in its
complement are displayed below.  Cycle labels are consecutive; a string
such as $03$ denotes the set $\{0,3\}$.
\[
\begin{array}{c|c|c}
F&\text{six branch sets}&\text{only missing pair}\\\hline
C_8(0,\ldots,7)&03,15,2,4,6,7&6:7\\
C_5(0,\ldots,4)\dot\cup C_3(5,6,7)&05,1,26,3,4,7&3:4\\
C_4(0,\ldots,3)\dot\cup C_4(4,\ldots,7)&04,1,25,3,6,7&6:7\\
C_7(0,\ldots,6)\dot\cup K_1(7)&024,1,3,5,6,7&5:6\\
C_4(0,\ldots,3)\dot\cup C_3(4,5,6)\dot\cup K_1(7)
 &024,1,3,5,6,7&5:6\\
C_6(0,\ldots,5)\dot\cup K_2(6,7)&024,1,3,5,6,7&6:7\\
C_3(0,1,2)\dot\cup C_3(3,4,5)\dot\cup K_2(6,7)
 &03,14,2,5,6,7&6:7
\end{array}
\]
Each nonsingleton is connected in $\overline F$, and all interbag
adjacencies except the indicated one are present.  Since $\overline F$
is a spanning subgraph of $J$, each displayed model also lies in $J$.
\end{proof}

Theorem~\ref{thm:defect-ladder} needs an edge of codegree at most three
without a colouring hypothesis.  Contracting exterior components reduces
the degree-seven and degree-eight cases to the following finite lemma.

\begin{lemma}[Computer-assisted finite lemma]\label{lem:finite-quotients}
For a graph $J$ and $A\subseteq V(J)$, let $Q(J,A)$ be obtained by adding
nonadjacent vertices $v,c$, making $v$ complete to $J$ and giving $c$
precisely the neighbourhood $A$.
\begin{enumerate}
\item If $|V(J)|=7$, $\delta(J)\geq4$ and $|A|\geq6$, then
$K_7^{-}\minor Q(J,A)$.
\item Fix a labelled graph $J$ with $|V(J)|=8$, $\delta(J)\geq4$ and
$K_6^{-}\not\minor J$.  At most one set $A\subseteq V(J)$ with $|A|\geq6$
makes $Q(J,A)$ $K_7^{-}$-minor-free.  If such a set exists,
$V(J)\setminus A$ consists of two adjacent vertices of degree four in $J$.
\end{enumerate}
\end{lemma}

\begin{proof}[Finite verification and its scope]
For part (1), $\overline J$ has maximum degree two, hence is a union of
paths and cycles.  Enumerating the multisets of such components with total
order seven gives 29 types.  For each, check the full attachment and its
seven possible single omissions.  Each of the $29\cdot8=232$ nine-vertex
quotients has a seven-bag model.  The checker enumerates all 750 partitions
of vertex subsets into seven nonempty bags, and tests connectivity and all
interbag adjacencies.  A seven-bag model on nine vertices uses seven, eight
or nine vertices, so this search is exhaustive.

For part (2), adjoin a vertex with each possible neighbourhood to each of
the 1,044 unlabelled seven-vertex graphs in the NetworkX graph atlas.
Every eight-vertex graph occurs.  The minimum-degree filter and exact
isomorphism testing leave 424 classes.
Of these, 55 have no $K_6^{-}$ minor.  Checking all missed sets of order
zero, one or two gives $55(1+8+28)=2,035$ exterior profiles.  Exactly 2,031
have certified $K_7^{-}$ models.  Each of the four remaining profiles
has the uniqueness and degree properties in part (2).

Here the minor engine starts with singleton bags and recursively either
deletes a bag or merges two touching bags.  These operations preserve
disjoint connected bags and generate every minor model by contracting
spanning trees and deleting unused vertices.  At six or seven bags it
accepts precisely when at most one pair is nonadjacent.  Every positive
certificate is separately checked for disjointness, connectivity and
interbag adjacencies.  The implementations test known positive and negative
instances; a separate internal audit also checked the four negative
quotients by all 11,880 seven-bag partitions of their vertex subsets.

The accompanying reproducibility note, \texttt{README.md}, records the
four profiles, executable inputs and SHA-256 certificate digests.
Part (1) uses only the Python standard library; part (2) also uses
NetworkX's atlas and exact isomorphism implementation.  The reductions
below impose no bound on the order of an exterior component or of $G$.
\end{proof}

\begin{theorem}\label{thm:six-connected-eight}
In every six-connected $K_7^{-}$-minor-free graph, each degree-eight vertex
is incident with an edge of codegree at most three.
\end{theorem}

\begin{proof}
Let $v$ have degree eight and put $J=G[N(v)]$.  Suppose every incident
codegree is at least four.  Then $\delta(J)\geq4$, and $J$ has no
$K_6^{-}$ minor because adjoining $\{v\}$ would give the forbidden
minor.  If $G-N[v]$ is empty, minimum degree six in $G$ gives
$\delta(J)\geq5$, contradicting Lemma~\ref{lem:eight-complement}.

Every component $C$ of $G-N[v]$ has at least six neighbours in $N(v)$,
since its neighbourhood separates it from $v$.  Contracting $C$ to $c$
and deleting the other exterior components gives exactly
$Q(J,N_G(C))$.  Lemma~\ref{lem:finite-quotients}(2) therefore applies to
every exterior component.  For the fixed labelled graph $J$, all
components must miss the same unique pair $a,b$.  Neither $a$ nor $b$
has an exterior neighbour, so both have degree $1+4=5$ in $G$,
contradicting six-connectivity.
\end{proof}

The degree-six case adapts a computation-free argument of Norin and
Totschnig.  We verify the hypotheses under $K_7^{-}$ exclusion.

\begin{lemma}[Degree-six disc bound]\label{lem:degree-six}
Let $H$ be a five-connected $K_7^{-}$-minor-free graph on at least nine
vertices.  If every edge has codegree at least four and some vertex has
degree six, then $|E(H)|\leq4|V(H)|-9$.
\end{lemma}

\begin{proof}
Let $v$ have degree six.  Since $d_{H[N(v)]}(x)=c(vx)\geq4$, label
$N(v)=\{u_1,w_1,u_2,w_2,u_3,w_3\}$ so that all its possible nonedges
are $u_iw_i$.  Two disjoint paths joining two distinct displayed pairs, with open
interiors outside $N[v]$,
would fill two of the three missing edges after contraction.  Together
with $v$ they give a $K_7^{-}$ model.  We call this the two-pair
obstruction; a present pair-edge may serve as one of the paths.

All three displayed pair-edges are absent.  Otherwise an exterior
component, with at least five neighbours in $N(v)$ by five-connectivity,
joins a second pair and violates the two-pair obstruction.  Moreover,
some pair, say $u_1,w_1$, has no exterior common neighbour.  Otherwise
choose a common neighbour for each pair.  Two distinct choices contradict
the obstruction.  If all choices coincide at $x$, order at least nine
leaves a component outside $N[v]\cup\{x\}$.  Its boundary has at least
five vertices, at least four of which belong to $N(v)$.  It therefore
joins one pair disjointly from a second pair's path through $x$, again a
contradiction.

We now use the two-paths theorem in the form of
\cite[Theorem~13]{NorinTotschnig2025}.  The next two applications are the
degree-six argument in \cite[Claims~3.12--3.15]{NorinTotschnig2025}; we
record the separator and density hypotheses needed for them.  There is
no nontrivial separation $(A,B)$ of order five with $N[v]\subseteq A$.
Indeed, five-connectivity and Menger's theorem provide five disjoint paths
in $H[A]$ from five distinct members of $N(v)$ to the five boundary
vertices, with interiors avoiding $N[v]$.  These five members contain
two complete matching pairs.  Temporarily label their boundary ends
$s_1,t_1$ and $s_2,t_2$, and label the fifth end $x$.  Apply the two-paths
theorem in $H[B]-x$ to the order $s_1,s_2,t_1,t_2$.  A crossing violates
the two-pair obstruction.  A separation outcome lifts after restoring
$x$ to a cut of order at most four in $H$.  In the remaining outcome the
graph embeds in a disc with the four terminals on its boundary.  An edge
with an end in $B\setminus A$ has at most two common neighbours in that
disc: three would enclose a vertex behind a triangle, which together
with $x$ gives a cut of order at most four in $H$.  Restoring $x$ gives
codegree at most three, a contradiction.

Apply the same theorem to $H'=H-\{v,u_1,w_1\}$, with terminal order
$u_2,u_3,w_2,w_3$.  A crossing again violates the two-pair obstruction.
A separating outcome lifts, after restoring $u_1,w_1$, to a separation
of order at most five with $N[v]$ on one side; five-connectivity and the
preceding paragraph exclude it.  Thus $H'$ embeds in a disc with the
four terminals on its boundary.  Put $n'=|V(H')|$.  Then
$|E(H')|\leq3n'-7$.  Each of its vertices outside $N[v]$ sees at most
one of $\{v,u_1,w_1\}$; the four remaining neighbourhood vertices see
all three, and the deleted triple induces two edges.  Hence
\[
 |E(H)|\leq(3n'-7)+(n'-4)+12+2=4n'+3=4|V(H)|-9.
\]
The forbidden-minor hypothesis was used only for the displayed
$K_7^{-}$ two-pair model.  Neither edge-maximality nor the stronger
$K_7^{\vee}$ exclusion is used in this argument.
\end{proof}

\begin{proposition}\label{prop:density-edge}
Let $G$ be six-connected and $K_7^{-}$-minor-free.  If
$|E(G)|\geq4|V(G)|$, then some edge of $G$ has codegree at most three.
\end{proposition}

\begin{proof}
The density implies order at least nine, so
Lemma~\ref{lem:jakobsen-defect} gives average degree less than nine.  Thus a
vertex of minimum degree has degree six, seven or eight.  Suppose every
edge has codegree at least four.  Lemma~\ref{lem:degree-six} excludes
degree six, and Theorem~\ref{thm:six-connected-eight} excludes degree
eight.  At a degree-seven vertex $v$, the local graph $J=G[N(v)]$ has
minimum degree at least four.  There is an exterior component, and
six-connectivity gives it at least six boundary neighbours.  Contracting
it produces a quotient forbidden by Lemma~\ref{lem:finite-quotients}(1).
This excludes all three degrees.
\end{proof}

\begin{proof}[Proof of Theorem~\ref{thm:defect-ladder}]
Let $G$ satisfy the hypotheses.  Proposition~\ref{prop:density-edge}
supplies an edge $e$ of codegree at most three.  The quotient $G/e$ is
$(r-1)$-connected and has no $K_7^{-}$ minor.  By \eqref{eq:contraction},
$|E(G/e)|\geq4|V(G/e)|$ and $D(G)\geq D(G/e)+1$.
The retained density ensures $|V(G/e)|\geq9$.

For $r=6$, Lemma~\ref{lem:jakobsen-defect} gives $D(G/e)\geq25$, so
$D(G)\geq26$.  For $r\geq7$, induction on $r$ applies to $G/e$ and gives
\[
              D(G)\geq D(G/e)+1\geq20+(r-1)+1=20+r.
\]
\end{proof}

The base case uses Jakobsen's bound because contraction need not preserve
six-connectivity.

\section{Seven-vertex separators}\label{sec:seven-cuts}

This section is computation-free.  A connected subgraph of $G-S$ is
\emph{full at $S$} if it is adjacent to every vertex of $S$.
We retain the exact equality partition of a boundary colouring, because
its mere number of colours does not suffice for gluing.

\begin{lemma}\label{lem:reflection}
Let $V(G)=L\mathbin{\dot\cup}S\mathbin{\dot\cup}R$, with $L,R$ nonempty
and anticomplete, and suppose every proper minor of $G$ is six-colourable.
Let $\Pi$ partition $S$ into independent blocks.  Suppose that all its
blocks, except some singleton blocks whose vertices form a clique, can
be assigned injectively to disjoint connected subgraphs of $G[L]$ full
at $S$, with at least one block assigned.  Then $G[R\cup S]$ has a
six-colouring with equality partition on $S$ exactly $\Pi$.
\end{lemma}

\begin{proof}
For each assigned block $B_j$ and its full subgraph $P_j$, contract the
connected set $B_j\cup V(P_j)$.  These sets are disjoint.  Their images,
together with the retained singleton blocks, form a clique: fullness
supplies all adjacencies involving an assigned block, and the retained
singletons form a clique by hypothesis.  At least one edge is contracted,
so the resulting minor is proper.  Six-colour it and pull back only to
$R\cup S$, assigning a block the colour of its image.  Each block is
independent and every edge to a retained vertex was represented in the
minor, so the pullback is proper.  The clique forces distinct colours on
distinct blocks, giving exactly $\Pi$.
\end{proof}

\begin{lemma}\label{lem:seven-small}
A graph on seven vertices with at least ten edges contains a $K_4^{-}$,
a house, or $K_{2,3}$ as a not necessarily induced subgraph.  A house is
a four-cycle with an additional vertex joined to the ends of one cycle
edge.
\end{lemma}

\begin{proof}
Suppose $Q$ contains none of the three.  Choose a maximum-degree vertex
$v$, put $d=d(v)$, $A=N(v)$, and $B=V(Q)\setminus N[v]$.
The graph $Q[A]$ is a matching, since a two-edge path there gives
$K_4^{-}$ with $v$.  Every vertex of $B$ has at most two neighbours in
$A$, since three would give $K_{2,3}$ with $v$.  For $d\geq5$,
\[
 |E(Q)|\leq d+\lfloor d/2\rfloor+2(6-d)+\binom{6-d}{2}\leq9.
\]
For $d\leq2$ there are at most seven edges.  Write
$\alpha=|E(Q[A])|$, $x=|E_Q(A,B)|$, $\beta=|E(Q[B])|$.

If $d=4$, then $\alpha\leq2$, $x\leq4$, $\beta\leq1$.
For $\alpha=2$, ten edges force a vertex of $B$ with two neighbours in
$A$.  If these are matched, there is a $K_4^{-}$; if they are on distinct
matching edges, there is a house.  For $\alpha\leq1$, ten edges force
$(\alpha,x,\beta)=(1,4,1)$.  Let $ab$ be the edge in $A$ and $c,t$ its
isolated vertices.  A two-set of neighbours in $A$ using $a$ or $b$
gives $K_4^{-}$ or a house.  Thus the two vertices $p,q$ of $B$ both see
$c,t$; their edge gives a $K_4^{-}$ on $p,q,c,t$.

If $d=3$, maximum degree gives exactly ten edges, so
$\alpha+x+\beta=7$ and $\alpha\leq1$.  For $\alpha=1$, no vertex of
$B$ can have two neighbours in $A$, by the same house or $K_4^{-}$
argument.  Hence $x=\beta=3$: $B$ is a triangle and its three edges to
$A$ have distinct ends, since a repetition gives $K_4^{-}$.  The edge
of $A$, its two matched vertices in $B$, and $v$ form a house.

It remains that $\alpha=0$.  If $\beta=3$, each vertex of $B$ already has two neighbours in $B$,
so $x\leq3$, a contradiction.  If $\beta=2$, write $B=pqr$ as a path.
Then $x=5$, and its cross-degrees are $2,1,2$.  Equal two-element
$A$-neighbourhoods of $p,r$ give $K_{2,3}$.  If they differ, the neighbour
of $q$ is either their common vertex, giving $K_4^{-}$, or one of the
two other vertices, giving a house.  If $\beta=1$, every vertex of $B$
has two neighbours in $A$.  Repeated two-sets give $K_4^{-}$ or
$K_{2,3}$.  Otherwise label the edge of $B$ as $pq$ and write
\[
 N_A(p)=\{a,b\},\quad N_A(q)=\{a,c\},\quad N_A(r)=\{b,c\}.
\]
Then $a,v,b,p$ is a four-cycle with roof $q$.  Finally $\beta=0$ would
force $x=7>6$.  Thus $Q$ has at most nine edges.
\end{proof}

\begin{lemma}\label{lem:three-cut-boundary}
Let $G$ satisfy the hypotheses of Theorem~\ref{thm:main}, and let $S$
be a seven-vertex cut.
Then $G-S$ has two or three components.  If it has three, $G[S]$ has a
proper three-colouring and every such colouring has class sizes $3,2,2$.
\end{lemma}

\begin{proof}
Every component is full at $S$ by seven-connectivity.  If there are at
least five, contract five of them to $c_1,\ldots,c_5$, label
$S=\{s_1,\ldots,s_7\}$, and use
\[
             \{c_i,s_i\}\ (1\leq i\leq5),\qquad\{s_6\},\{s_7\}.
\]
All but possibly the last pair are adjacent, giving $K_7^{-}$.

Suppose there are four components, contracted to $c_1,\ldots,c_4$.
If $G[S]$ has a path $xyz$, choose three further vertices $u,v,w$ of
$S$ and use
\[
 \{c_1\},\quad\{c_2,u\},\quad\{c_3,v\},\quad\{c_4,w\},
 \quad\{x\},\quad\{y\},\quad\{z\}.
\]
Only $\{x\}:\{z\}$ may be nonadjacent, again impossible.  Hence $G[S]$
is a matching and isolated vertices.  It has an isolated vertex $q$ and
a partition $S=I_1\mathbin{\dot\cup}I_2\mathbin{\dot\cup}\{q\}$ into
independent sets of orders $3,3,1$.  Split the four components into two
pairs, whose unions are $L,R$.  In each union the two full components
can be assigned to $I_1,I_2$ in Lemma~\ref{lem:reflection}, retaining
$q$.  The two colourings have the same exact boundary partition;
permuting colour names makes them agree on $S$ and yields a six-colouring
of $G$, a contradiction.  Thus there are at most three components.

Suppose there are three, with contracted vertices $c_1,c_2,c_3$.
If $G[S]$ had at least ten edges, Lemma~\ref{lem:seven-small} would
supply one of its three subgraphs.  For a $K_4^{-}$ subgraph, retain its
four vertices and merge each $c_i$ with one of the three unused boundary
vertices.  For a house or $K_{2,3}$, contract respectively the square
edge opposite the roof or any edge to obtain a four-bag $K_4^{-}$ model
using five boundary vertices.  Retain $c_1$ and merge $c_2,c_3$ with the
two unused boundary vertices.  In either case the seven sets form a
$K_7^{-}$ model.  Consequently $|E(G[S])|\leq9$.  The boundary also has
no $K_4$ subgraph, by the same construction using four singleton roots.

Every $K_4$-free graph on seven vertices with at most nine edges is
three-colourable.  Otherwise take a four-critical subgraph $F$; it has
minimum degree at least three.  On four vertices it is $K_4$.  On five
vertices parity gives a universal vertex, whose deletion leaves a
three-chromatic four-vertex graph, containing a triangle; this also gives
$K_4$.  On six vertices a degree at least four forces ten edges.  Otherwise
$F$ is cubic, with complement $C_6$ or $2C_3$; pairing consecutive
vertices in the former or taking the triangles in the latter gives a
three-colouring.  On seven vertices minimum degree three gives at least
eleven edges.  These contradictions prove the assertion.

If $G[S]$ were two-colourable, fix its partition into at most two
independent blocks.  For each component use the other two full components
in Lemma~\ref{lem:reflection} to reproduce the partition on that closed
component.  The three colourings align and glue, a contradiction.
Likewise, if any proper three-colouring has a singleton class, retain that
singleton and assign the other two blocks to the other two components.
The same gluing contradiction results.  Every class of a proper
three-colouring therefore has order at least two; their orders are
$3,2,2$.
\end{proof}

\begin{theorem}\label{thm:three-component-colouring}
Let $G$ be seven-connected, with every proper minor six-colourable.  Let
$S$ be a seven-vertex cut with exactly three components $D_1,D_2,D_3$.
If $G[S]$ has a proper colouring with class sizes $3,2,2$, then $G$ is
six-colourable.  No excluded-minor hypothesis is needed.
\end{theorem}

\begin{proof}
Fix the independent boundary blocks
$S=T\mathbin{\dot\cup}A\mathbin{\dot\cup}B$, with
$|T|=3$ and $|A|=|B|=2$.  Every $D_i$ is full at $S$.
Say that a component realises $(A,B)$ if it contains disjoint connected
subgraphs $X_A,X_B$, the first adjacent to both vertices of $A$, the
second to both vertices of $B$.  They can also be made adjacent: join
them by a shortest path in the component, split it at an edge, and absorb
the two halves into the respective subgraphs.

If two components realise $(A,B)$, then for every retained component
$D_i$ we may choose a different realising component $R$ and let $F$ be
the third.  Contract the three disjoint connected sets
\[
                X_A\cup A,\qquad X_B\cup B,\qquad F\cup T.
\]
Their images form a triangle, using the adjacency of $X_A$ and $X_B$
and fullness of $F$.  A six-colouring of this proper minor pulls back on
$G[D_i\cup S]$ with exact boundary partition $T|A|B$.  The three
colourings then align and glue.  Thus we may assume $D_1,D_2$ do not
realise $(A,B)$.

Fix one of these two components $D_i$.  Add four artificial terminals
$a_1^*,b_1^*,a_2^*,b_2^*$, each with the neighbours in $D_i$ of its
corresponding literal root.  There are no two disjoint paths joining the
$a$-pair and the $b$-pair, since their interiors would realise $(A,B)$.
Add the cycle edges in the displayed terminal order.  These do not
create such a crossing: a path using an edge between consecutive
terminals would meet a terminal of the other prescribed path.  Extend
on the same vertex set to an edge-maximal graph with no such crossing.
Humeau and Pous~\cite[Theorem~1.3]{HumeauPous2025} identify this completion
as a web with the four terminals as frame.  In this form a web consists
of a planar graph with the frame on a face and with cliques attached
behind facial triangles.

No clique inserted behind such a triangle can contain a vertex of $D_i$
outside the triangle.  Otherwise those original vertices have no
auxiliary neighbours outside the triangle.  Replace any artificial
terminal of the triangle by its literal root and add the three omitted
vertices of $T$.  These at most six actual vertices then separate a
nonempty part of $D_i$ from another component of $G-S$, contrary to
seven-connectivity.  Thus every original vertex of $D_i$ belongs to the
planar part of the web.  Replacing the artificial terminals by their
literal roots shows that
\[
                   H_i^+=G[D_i\cup A\cup B]+E(K_{A,B})
\]
is planar, with the added $K_{2,2}$ as an induced facial four-cycle.
The cycle is induced because $A$ and $B$ are independent.  These added
edges serve only in the planar extension; they are not asserted to be
edges of $G$ or used in a minor of $G$.

Contract in $G$ the two disjoint connected sets
\[
                        Q_0=D_1\cup T,\qquad Q_1=D_2\cup A.
\]
Their images are adjacent and each sees both retained vertices of $B$.
Six-colour this proper minor.  Pulling back on $G[D_3\cup S]$ gives a
colouring $c$ in which $T$ has colour $0$, $A$ a different colour $1$,
and both vertices of $B$ avoid $0,1$.

For $i=1,2$, precolour the induced four-cycle of $H_i^+$ by
$c|_{A\cup B}$.  It uses at most three colours from the five-colour
palette excluding $0$.  Diwan~\cite[Corollary~1]{Diwan2023} proves that
a proper precolouring of an induced cycle of length at most $2k-5$
using at most $k-1$ colours extends to a $k$-colouring of every planar
supergraph containing it.  With $k=5$ it extends to each $H_i^+$.
Use these extensions on $D_1,D_2$ and retain $c$ on $D_3\cup S$.
They agree on $A\cup B$ and avoid $0$ on $D_1,D_2$, so every edge to
$T$ is proper.  Distinct components are anticomplete.  The combined
colouring is a six-colouring of $G$.
\end{proof}


\begin{proof}[Proof of Theorem~\ref{thm:main}]
Proposition~\ref{prop:critical-base} gives seven-connectivity, minimum
degree eight, $|E(G)|\geq4|V(G)|$ and the absence of a $K_5$ subgraph.
The last conclusion makes every neighbourhood $K_4$-free.
Lemma~\ref{lem:independent-triple} gives independence number three in
every degree-eight neighbourhood; Proposition~\ref{prop:critical-three-edge}
gives the incident-edge conclusion and $n_8\geq26+\tau$ without computation.
Theorem~\ref{thm:defect-ladder} applied with $r=7$ gives
$9|V(G)|-2|E(G)|\geq27$.  The exact identity
$9|V(G)|-2|E(G)|=n_8-\tau$ proves the asserted stronger bound.
For a seven-vertex cut $S$, Lemma~\ref{lem:three-cut-boundary} leaves two
or three components.  Three would give a $3,2,2$ boundary colouring, and
Theorem~\ref{thm:three-component-colouring} would make $G$ six-colourable.
Thus $G-S$ has exactly two components.
\end{proof}

\section*{Statement on AI use and independence}

OpenAI models contributed to theorem discovery, proofs, verification and
drafting.  Cavit Erginsoy directed the work and made the final decisions.  Erginsoy is not affiliated with OpenAI, and OpenAI
neither sponsored nor verified the work.

{\hbadness=2000
\bibliographystyle{amsplain}
\bibliography{references}
}

\end{document}
